\chapter{Noisy Channels}\label{ch:channels}

\begin{center}\itshape
Noise does not make reliable communication impossible;\\
it only makes it slower, and by an amount one can compute.\\
\upshape(Shannon's theorem, in our words)
\end{center}

Before 1948, engineers knew one way to fight noise: repeat yourself. Send each bit three times and take a majority, and an error rate of one in ten becomes one in thirty-six; send it a hundred times and errors all but disappear, but so does the rate, which falls to one bit per hundred transmissions. Everyone took for granted that perfect reliability required the rate to go to zero.

Shannon's 1948 paper, in its part on ``the discrete channel with noise'', said that everyone was wrong. Every noisy channel has a number $C$, its \emph{capacity}, such that at any rate below $C$ one can communicate with an error probability as small as one wishes, while above $C$ it is impossible. A wire that flips each bit with probability $0.1$ has capacity $1-h(0.1)\approx0.531$: more than half a bit per use can be sent, essentially without error. The proof was as surprising as the statement. Shannon did not construct a good code; he chose the codewords \emph{at random} and showed that the average error over all codes is small, so that some code must be good. Paul Erd\H{o}s had used the same kind of argument the year before to bound Ramsey numbers (Tome~3, Chapter~12); in Shannon's hands it became the foundation of a branch of engineering.

It also left engineers with a riddle: if a random code is almost surely good, why is it so hard to find one that is good \emph{and} can be decoded in reasonable time? The riddle took forty-five years to answer. In 1993 Claude Berrou, Alain Glavieux and Punya Thitimajshima presented ``turbo codes'' that came within about half a decibel of Shannon's limit, and the community needed some time to believe their curves; we shall meet their cousins in Chapter~7. This chapter proves Shannon's theorem and its converse.

\begin{think}
A channel flips each bit with probability $p=0.1$. With the repetition code of length $3$ you send $\frac13$ bit per use, with error probability $0.028$. Do you think that one can send $\frac12$ bit per use with error probability below $10^{-9}$? What about $0.6$ bit per use?
\end{think}

\section{Channels and codes}\label{sec:dmc}

A \emph{discrete memoryless channel} has a finite input alphabet $A$, a finite output alphabet $B$ and a stochastic matrix $W(y\mid x)$, the probability of receiving $y$ when $x$ is sent. Used $n$ times, it acts independently on each use:
\[
W^n(y_1\cdots y_n\mid x_1\cdots x_n)=\prod_{i=1}^nW(y_i\mid x_i).
\]
The two fundamental examples (Figure~\ref{fig:channels}) are the \emph{binary symmetric channel} $\mathrm{BSC}(p)$, which flips each bit with probability $p$, and the \emph{binary erasure channel} $\mathrm{BEC}(\varepsilon)$, which delivers each bit intact or, with probability $\varepsilon$, replaces it by an erasure symbol ``?''; the receiver knows which bits it lost, but not their values.

\begin{figure}[ht]
\centering
\fig{ch03-channels}
\caption{Left: the binary symmetric channel and the binary erasure channel. Right: their capacities $1-h(p)$ and $1-\varepsilon$ (Proposition~\ref{prop:bscbec}); the dot marks $C=1-h(0.1)\approx0.531$. An erasure is cheaper than an error: $\mathrm{BEC}(2p)$ has more capacity than $\mathrm{BSC}(p)$ (Exercise~\ref{ex:erasureerror}).}
\label{fig:channels}
\end{figure}

\begin{definition}\label{def:code}
An $(M,n)$ \emph{code} consists of an encoder $f:\{1,\dots,M\}\to A^n$, whose values are the \emph{codewords}, and a decoder $g:B^n\to\{1,\dots,M\}$. Its \emph{rate} is $R=\frac{\log_2M}n$ bits per channel use. For a message $m$, the error probability is $\lambda_m=P\bigl(g(Y^n)\ne m\bigr)$ when $Y^n$ has law $W^n(\cdot\mid f(m))$; the \emph{average} and \emph{maximal} error probabilities are $P_e=\frac1M\sum_m\lambda_m$ and $\lambda_{\max}=\max_m\lambda_m$. A rate $R$ is \emph{achievable} if there are $(\lceil2^{nR}\rceil,n)$ codes with $\lambda_{\max}\to0$ as $n\to\infty$.
\end{definition}

The repetition code of length $n$ (odd), decoded by majority, has $M=2$, rate $\frac1n$ and, on $\mathrm{BSC}(p)$, error probability $P(\mathrm{Bin}(n,p)>\frac n2)$: $3p^2(1-p)+p^3=0.028$ for $n=3$ and $p=0.1$, about $8.6\times10^{-3}$ for $n=5$, and $e^{-\Theta(n)}$ in general, by Chernoff's bound. The rate goes to zero with the error. Shannon's theorem says this is a bad bargain.

\section{Capacity}\label{sec:capacity}

When the input is random with law $p$ on $A$, the pair (input, output) has law $p(x)W(y\mid x)$, and the mutual information $I(X;Y)$ measures how much the output tells about the input. It depends only on $p$ and $W$; write it $I(p,W)$.

\begin{definition}\label{def:capacity}
The \emph{capacity} of the channel $W$ is $C=\max_pI(p,W)$, the maximum over the laws $p$ on $A$.
\end{definition}

The maximum exists, because $I(p,W)$ is a continuous function of $p$ on the compact simplex of laws on $A$. It is also concave, which makes it easy to compute.

\begin{proposition}\label{prop:concave}
For fixed $W$, $p\mapsto I(p,W)$ is concave.
\end{proposition}

\begin{proof}
$I(p,W)=H(Y)-H(Y\mid X)$. The law of $Y$ is $q=pW$, a linear function of $p$, and $H$ is concave (Exercise~\ref{ex:convex} of Chapter~1), so $p\mapsto H(pW)$ is concave. And $H(Y\mid X)=\sum_xp(x)H\bigl(W(\cdot\mid x)\bigr)$ is linear in $p$.
\end{proof}

\begin{proposition}\label{prop:bscbec}
$\mathrm{BSC}(p)$ has capacity $1-h(p)$ and $\mathrm{BEC}(\varepsilon)$ has capacity $1-\varepsilon$; both are reached by the uniform input.
\end{proposition}

\begin{proof}
For the BSC, given $X=x$ the output is $x$ or its flip with probabilities $1-p$, $p$, so $H(Y\mid X)=h(p)$ whatever the law of $X$, and $I=H(Y)-h(p)\le1-h(p)$, with equality when $Y$ is uniform, which happens for the uniform input. For the BEC, compute the other way: $H(X\mid Y)$ is $0$ when $Y$ is not an erasure and $H(X)$ when it is, an event of probability $\varepsilon$ independent of $X$; so $I=H(X)-\varepsilon H(X)=(1-\varepsilon)H(X)\le1-\varepsilon$, with equality for the uniform input.
\end{proof}

For the BEC the capacity is obvious once seen: a fraction $1-\varepsilon$ of the bits arrive, and even a sender who knew in advance which positions would be erased could not send more. That one can reach $1-\varepsilon$ \emph{without} that knowledge is the content of Shannon's theorem for this channel. The value $C=1-h(0.1)\approx0.531$ answers the box: half a bit per use is possible, $0.6$ is not.

\section{Jointly typical pairs}\label{sec:jointtypical}

Fix a law $p$ on $A$ and let $(X_i,Y_i)$, $i=1,\dots,n$, be independent pairs with law $p(x)W(y\mid x)$; we may drop the letters of probability zero from $A$ and $B$. Write $p(x^n)$, $q(y^n)$ and $r(x^n,y^n)$ for the probabilities of sequences under the laws of $X^n$, $Y^n$ and $(X^n,Y^n)$.

\begin{definition}\label{def:jointtypical}
For $\varepsilon>0$, the set $J_\varepsilon^{(n)}$ of \emph{jointly typical} pairs consists of the $(x^n,y^n)$ such that
\[
\begin{gathered}
\Bigl|\frac1n\log_2\frac1{p(x^n)}-H(X)\Bigr|<\varepsilon,\qquad\Bigl|\frac1n\log_2\frac1{q(y^n)}-H(Y)\Bigr|<\varepsilon,\\
\Bigl|\frac1n\log_2\frac1{r(x^n,y^n)}-H(X,Y)\Bigr|<\varepsilon.
\end{gathered}
\]
\end{definition}

\begin{lemma}[joint typicality]\label{lem:jointtypical}
\begin{enumerate}[label=(\alph*)]
\item $P\bigl((X^n,Y^n)\in J_\varepsilon^{(n)}\bigr)\to1$ as $n\to\infty$.
\item $|J_\varepsilon^{(n)}|\le2^{n(H(X,Y)+\varepsilon)}$.
\item If $\tilde X^n$ and $\tilde Y^n$ are \emph{independent}, with the laws of $X^n$ and $Y^n$ respectively, then
\[
P\bigl((\tilde X^n,\tilde Y^n)\in J_\varepsilon^{(n)}\bigr)\le2^{-n(I(X;Y)-3\varepsilon)}.
\]
\end{enumerate}
\end{lemma}

\begin{proof}
(a) Each of the three conditions is the event that an average of $n$ independent, identically distributed, bounded variables ($\log_2\frac1{p(X_i)}$, $\log_2\frac1{q(Y_i)}$, $\log_2\frac1{r(X_i,Y_i)}$) is within $\varepsilon$ of its mean; by the weak law of large numbers each has probability tending to $1$, and so does their intersection. (b) Every element of $J_\varepsilon^{(n)}$ has probability $r(x^n,y^n)\ge2^{-n(H(X,Y)+\varepsilon)}$, and these probabilities add up to at most $1$. (c) The pair $(\tilde X^n,\tilde Y^n)$ takes the value $(x^n,y^n)$ with probability $p(x^n)q(y^n)$, which for a jointly typical pair is at most $2^{-n(H(X)-\varepsilon)}2^{-n(H(Y)-\varepsilon)}$. Summing over $J_\varepsilon^{(n)}$ and using (b) gives at most $2^{n(H(X,Y)+\varepsilon-H(X)-H(Y)+2\varepsilon)}=2^{-n(I(X;Y)-3\varepsilon)}$.
\end{proof}

Here is the picture behind the lemma. There are about $2^{nH(Y)}$ typical outputs. Sending a typical $x^n$ produces one of about $2^{nH(Y\mid X)}$ outputs that are jointly typical with it, a ``noise cloud''. If the clouds of the codewords are to be nearly disjoint, there can be at most $2^{nH(Y)}/2^{nH(Y\mid X)}=2^{nI(X;Y)}$ codewords. Shannon's theorem says that this count is right in both directions.

\section{Shannon's theorem}\label{sec:shannon}

\begin{theorem}[noisy channel coding theorem, Shannon 1948]\label{thm:shannon}
Every rate $R<C$ is achievable: there are $(\lceil2^{nR}\rceil,n)$ codes whose maximal error probability tends to $0$.
\end{theorem}

\begin{proof}
Let $p$ be a law on $A$ with $I(p,W)=C$, let $0<R<C$, choose $\varepsilon>0$ with $R+3\varepsilon<C$ and let $M=2\lceil2^{nR}\rceil$.

\emph{Random code.} Choose the $M$ codewords $c_1,\dots,c_M\in A^n$ independently, each with independent letters of law $p$. \emph{Decoder.} On receiving $y^n$, output the unique $m$ such that $(c_m,y^n)\in J_\varepsilon^{(n)}$; if there is none, or more than one, output $1$ (an error, unless $m=1$ was sent, which can only help).

\emph{Error, averaged over codes.} Let $E$ be the event that the decoder errs when message $m$ is sent, the probability being over the random code and the channel. By symmetry it does not depend on $m$; take $m=1$. The decoder can err only if $(c_1,Y^n)\notin J_\varepsilon^{(n)}$, or if $(c_k,Y^n)\in J_\varepsilon^{(n)}$ for some $k\ne1$. The pair $(c_1,Y^n)$ has the law of $(X^n,Y^n)$ of the lemma, so the first event has probability $\eta_n\to0$ by (a). For $k\ne1$, the codeword $c_k$ is independent of $c_1$ and of the noise, hence of $Y^n$, and has the law of $X^n$, while $Y^n$ has the law of the lemma's $Y^n$; by (c), $(c_k,Y^n)$ is jointly typical with probability at most $2^{-n(C-3\varepsilon)}$. By the union bound,
\[
P(E)\le\eta_n+(M-1)2^{-n(C-3\varepsilon)}\le\eta_n+4\cdot2^{-n(C-3\varepsilon-R)}\longrightarrow0,
\]
as $M\le2\cdot2^{nR}+2\le4\cdot2^{nR}$ (for $2^{nR}\ge1$).

\emph{From average to maximal error.} $P(E)$ is the average over random codes of the average error $P_e$ of the code. So some code $\mathcal C$ has $P_e(\mathcal C)\le P(E)=:\delta_n$. At most half of its messages have $\lambda_m>2\delta_n$ (otherwise the average would exceed $\delta_n$). Remove them, keeping the same decoder (restricted to the remaining messages, it can only err less): the $\frac M2=\lceil2^{nR}\rceil$ remaining codewords form a code with $\lambda_{\max}\le2\delta_n\to0$.
\end{proof}

The proof is an existence proof twice over: it does not say which code works, and the jointly typical decoder looks through all $M$ codewords. A random code needs an exponential table even to be described. Linear codes repair the first defect at no cost, at least on the binary symmetric channel.

\begin{theorem}[Elias 1955]\label{thm:linearcap}
Let $0<p<\frac12$ and $R<1-h(p)$. Let $G$ be a uniformly random $k\times n$ matrix over $\mathbb F_2$ with $k=\lceil Rn\rceil$, encode $m\in\mathbb F_2^k$ as $mG$, and decode $y$ to a message whose codeword is nearest to $y$ in Hamming distance. On $\mathrm{BSC}(p)$, the error probability, averaged over $G$, tends to $0$ for every message. In particular there are linear codes of rate at least $R$ with average error probability tending to $0$.
\end{theorem}

\begin{proof}
Let $\delta>0$ with $p+\delta\le\frac12$ and $R<1-h(p+\delta)$, and $r=n(p+\delta)$. The noise is a word $e$ with independent bits of mean $p$, independent of $G$; the message $m$ is sent and $y=mG+e$ received. The decoder can err only if $e$ has weight more than $r$, or some $m'\ne m$ has $d(m'G,y)\le r$. The first event has probability tending to $0$ by the weak law. For the second, $y-m'G=e+(m-m')G$, and $(m-m')G$ is a uniformly random word of $\mathbb F_2^n$, independent of $e$: indeed, a fixed nonzero $u\in\mathbb F_2^k$ has a nonzero coordinate $u_i$, and conditionally on the other rows of $G$, row $i$ is uniform, so $uG$ is uniform. Hence $y-m'G$ is uniform, and by Proposition~\ref{prop:binomial},
\[
P\bigl(d(m'G,y)\le r\bigr)=2^{-n}\sum_{j\le r}\binom nj\le2^{-n(1-h(p+\delta))}.
\]
The union bound over the $2^k-1$ messages $m'\ne m$ gives at most $2^{k-n(1-h(p+\delta))}\le2\cdot2^{-n(1-h(p+\delta)-R)}\to0$. For the last claim, average over messages and choose a good $G$; its rank is $k$ for large $n$, since otherwise two messages would have the same codeword and the error would be at least $\frac12$ for half of them.
\end{proof}

So linear codes, which can be described by $kn$ bits and encoded by a matrix product, achieve capacity; decoding them remains hard (finding the nearest codeword in a linear code is NP-hard, a theorem of Elwyn Berlekamp, Robert McEliece and Henk van Tilborg from 1978, quoted). The codes of Chapters~5 to~7 trade a little rate for efficient decoding.

\section{The converse}\label{sec:converse}

\begin{theorem}[converse]\label{thm:converse}
Every $(M,n)$ code with $M\ge2$ and rate $R=\frac{\log_2M}n$ has average error probability
\[
P_e\ge1-\frac CR-\frac1{nR}.
\]
In particular no rate $R>C$ is achievable, even with $P_e$ in place of $\lambda_{\max}$.
\end{theorem}

\begin{proof}
Let the message $V$ be uniform on $\{1,\dots,M\}$, $X^n=f(V)$, $Y^n$ the output and $\hat V=g(Y^n)$, so that $P(\hat V\ne V)=P_e$. Then $nR=\log_2M=H(V)=H(V\mid Y^n)+I(V;Y^n)$. By Fano's inequality (Theorem~\ref{thm:fano}), $H(V\mid Y^n)\le1+P_e\log_2M=1+P_enR$. Since $V\to X^n\to Y^n$ is a Markov chain, $I(V;Y^n)\le I(X^n;Y^n)$ (Theorem~\ref{thm:dataprocessing}). Finally, the channel is memoryless, so $\log_2\frac1{W^n(y^n\mid x^n)}=\sum_i\log_2\frac1{W(y_i\mid x_i)}$ and, taking expectations, $H(Y^n\mid X^n)=\sum_iH(Y_i\mid X_i)$; with subadditivity,
\[
\begin{aligned}I(X^n;Y^n)&=H(Y^n)-\sum_iH(Y_i\mid X_i)\eqbr\le\sum_i\bigl(H(Y_i)-H(Y_i\mid X_i)\bigr)=\sum_iI(X_i;Y_i)\le nC.\end{aligned}
\]
Altogether $nR\le1+P_enR+nC$, which is the claim.
\end{proof}

For the binary symmetric channel a counting argument gives much more: above capacity, the probability of decoding correctly does not merely stay below $1$, it tends to $0$.

\begin{theorem}[strong converse for the BSC]\label{thm:strongconverse}
Let $0<p<\frac12$ and $R>1-h(p)$. Every code of length $n$ for $\mathrm{BSC}(p)$ with at least $2^{nR}$ codewords has average probability of correct decoding at most $\eta_n$, where $\eta_n\to0$ does not depend on the code.
\end{theorem}

\begin{proof}
Let $D_m=g^{-1}(m)$ be the set of received words decoded to $m$; these sets are disjoint, so $\sum_m|D_m|\le2^n$. Choose $\delta>0$ with $\delta\log_2\frac{1-p}p<R-1+h(p)$, and call the noise \emph{normal} if its weight $d$ satisfies $|d-np|\le\delta n$; abnormal noise has probability $\alpha_n\to0$ by the weak law. A given received word at distance $d$ from the codeword sent has probability $p^d(1-p)^{n-d}$, which decreases with $d$ since $p<\frac12$; for normal noise it is at most
\[
\begin{aligned}p^{n(p-\delta)}(1-p)^{n(1-p+\delta)}&=2^{-nh(p)}\Bigl(\frac{1-p}p\Bigr)^{n\delta}\eqbr=2^{-n(h(p)-\delta\log_2\frac{1-p}p)}.\end{aligned}
\]
So the probability that message $m$ is decoded correctly is at most $\alpha_n+|D_m|\,2^{-n(h(p)-\delta\log_2\frac{1-p}p)}$, and averaging over the $M\ge2^{nR}$ messages,
\[
\begin{aligned}\frac1M\sum_m(1-\lambda_m)&\le\alpha_n+\frac{2^n}{2^{nR}}2^{-n(h(p)-\delta\log_2\frac{1-p}p)}\eqbr=\alpha_n+2^{-n(R-1+h(p)-\delta\log_2\frac{1-p}p)}\longrightarrow0.\end{aligned}\qedhere
\]
\end{proof}

Capacity is a sharp threshold, like entropy in Tome~3, Corollary~4.22: below it errors can be made as rare as we like, above it decoding fails almost surely.

\section{Computing capacity}\label{sec:blahut}

For symmetric channels the uniform input is optimal (Exercise~\ref{ex:symmetric}). In general the maximisation is a concave problem in $|A|-1$ variables, and Suguru Arimoto and Richard Blahut found in 1972 an elegant algorithm for it, based on a second description of mutual information.

\begin{proposition}\label{prop:variational}
For a law $p$ on $A$ and any family of laws $Q(\cdot\mid y)$ on $A$, indexed by $y\in B$,
\[
I(p,W)\ge F(p,Q)=\sum_{x,y}p(x)W(y\mid x)\log_2\frac{Q(x\mid y)}{p(x)},
\]
with equality when $Q(x\mid y)$ is the posterior law $P(X=x\mid Y=y)$. For fixed $Q$, $F(p,Q)$ is maximal for $p(x)$ proportional to $\prod_y Q(x\mid y)^{W(y\mid x)}$.
\end{proposition}

\begin{proof}
Let $\pi(x\mid y)=P(X=x\mid Y=y)$ and $q(y)=P(Y=y)$. Since $p(x)W(y\mid x)=q(y)\pi(x\mid y)$,
\[
\begin{aligned}I(p,W)-F(p,Q)&=\sum_{x,y}q(y)\pi(x\mid y)\log_2\frac{\pi(x\mid y)}{Q(x\mid y)}\eqbr=\sum_yq(y)\,D\bigl(\pi(\cdot\mid y)\,\big\|\,Q(\cdot\mid y)\bigr)\ge0,\end{aligned}
\]
with equality for $Q=\pi$. For fixed $Q$, let $s(x)=\prod_yQ(x\mid y)^{W(y\mid x)}$ and $Z=\sum_xs(x)$. Then $F(p,Q)=\sum_xp(x)\log_2\frac{s(x)}{p(x)}=\log_2Z-D(p\,\|\,s/Z)$, maximal for $p=s/Z$.
\end{proof}

So $C=\max_p\max_QF(p,Q)$, and the \emph{Blahut--Arimoto algorithm} maximises alternately: from $p_t$, set $Q_t$ to the posterior of $p_t$, then $p_{t+1}\propto\prod_yQ_t(x\mid y)^{W(y\mid x)}$. Each step can only increase $F$, and $F(p_{t+1},Q_t)\le I(p_{t+1},W)\le C$. Arimoto and Blahut proved that, from any starting law with full support, $I(p_t,W)\to C$; started from the uniform law, the error after $t$ steps is at most $\frac1t\log_2|A|$ (quoted). The notebook runs it on the Z-channel of Exercise~\ref{ex:zchannel}.

\begin{remark}[continuous channels]
Shannon's paper also treats a channel with real inputs of power at most $P$ and additive Gaussian noise of power $N$; its capacity is $\frac12\log_2(1+\frac PN)$ bits per use, and for a band of $w$ hertz, which carries $2w$ independent samples per second, $w\log_2(1+\frac PN)$ bits per second (quoted; the proof follows this chapter with densities in place of laws). For a telephone line, with $w\approx3000$ hertz and a signal-to-noise ratio $P/N$ of about $35$ decibels, this is about $35\,000$ bits per second, which is where the fastest analogue modems of the 1990s stopped.
\end{remark}

\begin{beauty}
Shannon's theorem turns an engineering question into a single number, and its proof turns randomness into a tool: we cannot exhibit a good code, but we can prove that almost every code is good, by counting noise clouds and typical outputs. The converse is the same counting read backwards. Between them, the jointly typical set does all the work, and the AEP of Tome~3, a consequence of the weak law of large numbers, carries the whole theory.
\end{beauty}

\section{Exercises}

\begin{exercise}[$\star$]\label{ex:repetition}
Compute the error probability of the repetition codes of length $3$, $5$ and $7$ on $\mathrm{BSC}(0.1)$, and the length needed to bring it below $10^{-9}$. What is then the rate?
\end{exercise}

\begin{exercise}[$\star$]\label{ex:erasureerror}
Show that $h(p)\ge2p$ for $0\le p\le\frac12$, so that $\mathrm{BEC}(2p)$ has at least the capacity of $\mathrm{BSC}(p)$.
\end{exercise}

\begin{exercise}[$\star\star$]\label{ex:zchannel}
The \emph{Z-channel} transmits $0$ correctly and turns $1$ into $0$ with probability $\frac12$. Show that its capacity is $\log_2\frac54\approx0.322$, reached when $P(X=1)=\frac25$.
\end{exercise}

\begin{exercise}[$\star\star$]\label{ex:symmetric}
A channel is \emph{symmetric} if the rows of $W$ are permutations of each other, and so are the columns. Show that its capacity is $\log_2|B|-H(\text{row})$, reached by the uniform input.
\end{exercise}

\begin{exercise}[$\star$]\label{ex:cascade}
Two binary symmetric channels $\mathrm{BSC}(p)$ and $\mathrm{BSC}(q)$ in a row form $\mathrm{BSC}(p+q-2pq)$. Check that its capacity is at most the smaller of the two, as data processing predicts.
\end{exercise}

\begin{exercise}[$\star\star$]\label{ex:parallel}
Two channels $W_1$, $W_2$ used side by side form a channel with input alphabet $A_1\times A_2$. Show that its capacity is $C_1+C_2$.
\end{exercise}

\begin{exercise}[$\star$]\label{ex:fanobound}
Use Theorem~\ref{thm:converse} to show that a code of rate $0.6$ on $\mathrm{BSC}(0.1)$ has average error at least about $0.115$ for large $n$, and compute the bound for $n=100$.
\end{exercise}

\begin{exercise}[$\star\star$]\label{ex:spherepacking}
A code $\mathcal C\subseteq\{0,1\}^n$ corrects all patterns of at most $t$ errors if the Hamming balls of radius $t$ around the codewords are disjoint. Show that then $|\mathcal C|\sum_{j\le t}\binom nj\le2^n$, and deduce that if $t=\lfloor\tau n\rfloor$ with $\tau<\frac12$ fixed, the rate is at most $1-h(\tau)+o(1)$.
\end{exercise}

\begin{exercise}[$\star\star$]\label{ex:zerorate}
Show that no code with two or more codewords can correct \emph{every} pattern of at most $\lceil\frac n2\rceil$ errors, although for $p$ just below $\frac12$ Shannon's theorem gives codes of positive rate $1-h(p)-o(1)$ whose error probability on $\mathrm{BSC}(p)$, where about $pn$ bits are flipped, tends to $0$. Explain why there is no contradiction.
\end{exercise}

\begin{exercise}[$\star\star\star$]\label{ex:convexW}
Show that for fixed $p$, $I(p,W)$ is a convex function of the channel $W$.
\end{exercise}

\section{Hints}

\begin{hint}{ex:repetition}
The error is $P(\mathrm{Bin}(n,0.1)>\frac n2)$.
\end{hint}
\begin{hint}{ex:erasureerror}
$h$ is concave and $h(0)=0$, $h(\frac12)=1$.
\end{hint}
\begin{hint}{ex:zchannel}
With $a=P(X=1)$, $I=h(\frac a2)-a$.
\end{hint}
\begin{hint}{ex:symmetric}
$H(Y\mid X)$ is the entropy of a row for every input law.
\end{hint}
\begin{hint}{ex:cascade}
Compute the probability that exactly one of the two channels flips.
\end{hint}
\begin{hint}{ex:parallel}
Independent inputs give $\ge$; subadditivity gives $\le$.
\end{hint}
\begin{hint}{ex:fanobound}
$1-\frac CR$ with $C\approx0.531$.
\end{hint}
\begin{hint}{ex:spherepacking}
Count the words in the disjoint balls; Proposition~\ref{prop:binomial}.
\end{hint}
\begin{hint}{ex:zerorate}
For any two codewords, some pattern of at most $\lceil\frac n2\rceil$ errors turns one into a word at least as close to the other.
\end{hint}
\begin{hint}{ex:convexW}
Write $I(p,W)=D\bigl(p\otimes W\,\|\,p\otimes(pW)\bigr)$ and use the joint convexity of $D$.
\end{hint}

\section{Solutions}

\begin{solution}{ex:repetition}
\idea Majority decoding fails when more than half the bits flip: a binomial tail.
\writeup For $n=3$: $3(0.1)^2(0.9)+(0.1)^3=0.028$. For $n=5$: $\binom53(0.1)^3(0.9)^2+\binom54(0.1)^4(0.9)+(0.1)^5\approx0.00856$. For $n=7$: about $0.00273$. The error first drops below $10^{-9}$ at $n=35$ (notebook, among odd $n$; $n=33$ gives $2.4\times10^{-9}$), with rate $\frac1{35}\approx0.029$, while Shannon promises rates up to $0.531$ at any reliability.
\end{solution}

\begin{solution}{ex:erasureerror}
\idea A concave function lies above its chords.
\writeup On $[0,\frac12]$, $h$ is concave with $h(0)=0$ and $h(\frac12)=1$, so it lies above the chord $t\mapsto2t$: $h(p)\ge2p$. Hence $1-2p\ge1-h(p)$. An error costs as much as at least two erasures: the receiver of an erasure knows where the damage is.
\end{solution}

\begin{solution}{ex:zchannel}
\idea Compute $I=H(Y)-H(Y\mid X)$ as a function of the single parameter $a=P(X=1)$ and maximise.
\writeup $P(Y=1)=\frac a2$, so $H(Y)=h(\frac a2)$; given $X=0$ the output is certain, given $X=1$ it is a fair coin, so $H(Y\mid X)=a$. Thus $I(a)=h(\frac a2)-a$, a concave function with $I'(a)=\frac12\log_2\frac{1-a/2}{a/2}-1$. It vanishes when $\frac{1-a/2}{a/2}=4$, that is $a=\frac25$. Then $I=h(\frac15)-\frac25=\frac15\log_25+\frac45\log_2\frac54-\frac25=\log_25-\frac85-\frac25=\log_25-2=\log_2\frac54\approx0.322$. The optimal input sends $1$ less often than $0$, because $1$ is the fragile symbol.
\end{solution}

\begin{solution}{ex:symmetric}
\idea Symmetry fixes $H(Y\mid X)$, and makes the uniform input produce a uniform output.
\writeup Since each row of $W$ is a permutation of the first, $H(Y\mid X=x)$ is the entropy $H(\text{row})$ of a row, for every $x$; so $I=H(Y)-H(\text{row})\le\log_2|B|-H(\text{row})$. For the uniform input, $P(Y=y)=\frac1{|A|}\sum_xW(y\mid x)$, the sum of column $y$ divided by $|A|$; all columns are permutations of each other, so they have the same sum, and $Y$ is uniform. So the bound is reached.
\end{solution}

\begin{solution}{ex:cascade}
\idea The bit arrives flipped if exactly one of the two channels flips it.
\writeup The cascade flips with probability $p(1-q)+(1-p)q=p+q-2pq$. If $p,q\le\frac12$, then $\frac12-(p+q-2pq)=2(\frac12-p)(\frac12-q)$, which is at most both $\frac12-p$ and $\frac12-q$; as $h$ is increasing on $[0,\frac12]$, the cascade's capacity $1-h(p+q-2pq)$ is at most both capacities. (If $p>\frac12$, replace it by $1-p$ by relabelling the output, which does not change capacity.) This is the data processing inequality $I(X;Z)\le I(X;Y)$ for $X\to Y\to Z$, and also $I(X;Z)\le I(Y;Z)$.
\end{solution}

\begin{solution}{ex:parallel}
\idea Use the two channels independently to get at least $C_1+C_2$; subadditivity shows one cannot do better.
\writeup With independent inputs $X_1$, $X_2$ of optimal laws, the outputs are independent and $I(X_1X_2;Y_1Y_2)=I(X_1;Y_1)+I(X_2;Y_2)=C_1+C_2$, by the additivity of entropy for independent pairs. Conversely, for any joint input law, $H(Y_1Y_2\mid X_1X_2)=H(Y_1\mid X_1)+H(Y_2\mid X_2)$ because the channels act independently, and $H(Y_1Y_2)\le H(Y_1)+H(Y_2)$, so $I\le I(X_1;Y_1)+I(X_2;Y_2)\le C_1+C_2$.
\end{solution}

\begin{solution}{ex:fanobound}
\idea Plug numbers into $P_e\ge1-\frac CR-\frac1{nR}$.
\writeup With $C=1-h(0.1)\approx0.5310$ and $R=0.6$, $1-\frac CR\approx0.115$. For $n=100$ the correction $\frac1{nR}\approx0.0167$ leaves $P_e\ge0.098$. The strong converse says much more: the error tends to $1$.
\end{solution}

\begin{solution}{ex:spherepacking}
\idea Disjoint balls of the same volume in a space of $2^n$ points: count.
\writeup The balls of radius $t$ around the codewords are disjoint subsets of $\{0,1\}^n$ and each has $\sum_{j\le t}\binom nj$ elements, so $|\mathcal C|\sum_{j\le t}\binom nj\le2^n$. By Proposition~\ref{prop:binomial}, $\sum_{j\le t}\binom nj\ge\binom nt\ge\frac1{n+1}2^{nh(t/n)}$, so $\frac1n\log_2|\mathcal C|\le1-h(\frac tn)+\frac{\log_2(n+1)}n$, and $h(\frac{\lfloor\tau n\rfloor}n)\to h(\tau)$. This is the Hamming bound of Chapter~5.
\end{solution}

\begin{solution}{ex:zerorate}
\idea Average error and worst-case error are different requirements: Shannon allows a small fraction of noise patterns to defeat the decoder.
\writeup The positive statement is Theorem~\ref{thm:linearcap} (or Theorem~\ref{thm:shannon}). For the negative one, let $c\ne c'$ be codewords at distance $d\le n$, and let $e$ flip $\lceil\frac d2\rceil$ of the positions where they differ. The received word $y=c+e$ differs from $c$ in $\lceil\frac d2\rceil$ places and from $c'$ in $\lfloor\frac d2\rfloor$ places, so $y=c'+e'$ with $\mathrm{wt}(e')=\lfloor\frac d2\rfloor$. Both error patterns have weight at most $\lceil\frac n2\rceil$, so a decoder correcting all of them would have to decode $y$ both to $c$ and to $c'$: impossible.

No contradiction with Shannon: correcting \emph{all} patterns of $pn$ errors is a worst-case requirement, while on $\mathrm{BSC}(p)$ the noise is random, and the decoder is allowed to fail on a set of noise patterns of vanishing probability. Two codewords at distance $d$ are confused only when the noise flips at least half of the $d$ positions where they differ; random noise flips about $pd<\frac d2$ of them, so when $d$ grows linearly with $n$ this is a large deviation, of exponentially small probability. The worst-case problem has its own, much smaller, limits (Chapter~5).
\end{solution}

\begin{solution}{ex:convexW}
\idea Mutual information is a divergence between the joint law and the product of the marginals, and both depend linearly on $W$ for fixed $p$.
\writeup By Theorem~\ref{thm:mutualpos}, $I(p,W)=D\bigl(P_W\,\|\,Q_W\bigr)$ where $P_W(x,y)=p(x)W(y\mid x)$ and $Q_W(x,y)=p(x)\,(pW)(y)$. Both $P_W$ and $Q_W$ are linear in $W$. If $W=\lambda W_1+(1-\lambda)W_2$, then $P_W=\lambda P_{W_1}+(1-\lambda)P_{W_2}$ and similarly for $Q$, and the joint convexity of $D$ (Exercise~\ref{ex:convex} of Chapter~1) gives $I(p,W)\le\lambda I(p,W_1)+(1-\lambda)I(p,W_2)$. A mixture of channels is noisier than the average of its components, because the receiver does not know which component acted.
\end{solution}

\section*{Sources and further reading}

\begin{itemize}
\item \emph{A research paper to read.} C.~E.~Shannon, \emph{A mathematical theory of communication}, Bell System Technical Journal 27 (1948), 379--423 and 623--656, Part~II (Sections~13--19). Section~13 introduces the discrete channel with noise and the equivocation $H(X\mid Y)$, with the example of a channel that errs once in a hundred and the observation that its capacity is not $0.99$ but $0.919$ bit per symbol; Theorem~11 is our Theorem~\ref{thm:shannon}, proved in Section~14 by the random coding argument with typical sequences (compare with our Lemma~\ref{lem:jointtypical}); Section~15 computes capacities of symmetric channels. Then read Part~IV for the Gaussian channel of our remark. Shannon's proofs are sketches by modern standards; making them rigorous took a decade (Feinstein 1954, Wolfowitz 1957).
\item \emph{The stories.} C.~Berrou, A.~Glavieux and P.~Thitimajshima, \emph{Near Shannon limit error-correcting coding and decoding: turbo-codes}, Proceedings of the IEEE International Conference on Communications, Geneva, 1993, 1064--1070. P.~Erd\H{o}s, \emph{Some remarks on the theory of graphs}, Bulletin of the American Mathematical Society 53 (1947), 292--294.
\item P.~Elias, \emph{Coding for noisy channels}, IRE Convention Record 3, part~4 (1955), 37--46; A.~Feinstein, \emph{A new basic theorem of information theory}, IRE Transactions on Information Theory 4 (1954), 2--22; J.~Wolfowitz, \emph{The coding of messages subject to chance errors}, Illinois Journal of Mathematics 1 (1957), 591--606 (the strong converse); E.~R.~Berlekamp, R.~J.~McEliece and H.~C.~A.~van Tilborg, \emph{On the inherent intractability of certain coding problems}, IEEE Transactions on Information Theory 24 (1978), 384--386.
\item S.~Arimoto, \emph{An algorithm for computing the capacity of arbitrary discrete memoryless channels}, IEEE Transactions on Information Theory 18 (1972), 14--20; R.~E.~Blahut, \emph{Computation of channel capacity and rate-distortion functions}, IEEE Transactions on Information Theory 18 (1972), 460--473.
\item T.~M.~Cover and J.~A.~Thomas, \emph{Elements of Information Theory}, 2nd ed., Wiley, 2006, Chapters~7 and~9; R.~G.~Gallager, \emph{Information Theory and Reliable Communication}, Wiley, 1968, for error exponents; Y.~Polyanskiy and Y.~Wu, \emph{Information Theory: From Coding to Learning}, Cambridge University Press, 2025, a modern treatment with finite-length bounds.
\end{itemize}
